\chapter{RESULTADOS}

\section{Resultados}

Los resultados obtenidos durante el desarrollo  del proyecto LIA fueron evaluados en función de los objetivos planteados. A continuación, se presentan los hallazgos más relevantes, organizados en categorías clave:

\subsection{Implementación del Framework IoT}
El framework IoT diseñado permitió la integración de dispositivos de monitoreo ambiental en un entorno controlado. Los principales logros incluyen:
\begin{itemize}
    \item \textbf{Interfaz de usuario intuitiva:} Se desarrolló una plataforma interactiva que permite graficar y agrupar datos por fechas, descargar gráficos en formatos estándar y realizar ajustes en tiempo real.
    \item \textbf{Modularidad y escalabilidad:} El framework puede integrarse con nuevos dispositivos y adaptarse a diferentes especies de plantas, lo que asegura su aplicabilidad en otros proyectos agrícolas.
\end{itemize}

\subsection{Visualización y Análisis de Datos}
Se implementaron herramientas avanzadas para analizar y visualizar los datos recolectados:
\begin{itemize}
    \item \textbf{Modelos gráficos:} Se generaron gráficos dinámicos que muestran tendencias y patrones, permitiendo identificar correlaciones entre variables ambientales y el crecimiento de las plantas.
    \item \textbf{Descarga de reportes:} Los gráficos pueden descargarse en formatos estándar (\textit{e.g.}, PNG), facilitando el análisis externo y la presentación de resultados.
    \item \textbf{Alertas configurables:} El sistema incluye un módulo para generar alertas personalizadas basadas en umbrales definidos por el usuario.
\end{itemize}


\subsection{Conclusión de los Resultados}
El proyecto IOpenDataLink logró cumplir con los objetivos planteados, proporcionando una solución tecnológica innovadora para la agricultura moderna. El framework IoT desarrollado facilita el monitoreo, análisis y control de variables ambientales, mejorando la eficiencia y sostenibilidad de los cultivos en invernaderos. Estos resultados abren la posibilidad de extender la solución a entornos agrícolas más amplios y explorar nuevas aplicaciones de la tecnología IoT en la agricultura.
